% concatenated from main.tex
\documentclass[11pt]{amsart}
\usepackage{times}
\usepackage[T1]{fontenc}
\usepackage{amssymb, amsthm, amsmath}
\usepackage{mathrsfs}
\usepackage{mathtools}

\usepackage[active]{srcltx}

%\usepackage{dsfont}
%\usepackage[dvipdf]{graphicx}
\usepackage{todonotes}

\usepackage{setspace}
%\usepackage[margin=.8in]{geometry}
\usepackage{geometry}

\usepackage{hyperref}
%\usepackage{showlabels}

%\DeclareMathOperator{\cl}{cl}
\DeclareMathOperator{\dM}{DM}
\DeclareMathOperator{\stcl}{stcl}
\DeclareMathOperator{\acl}{acl}
\DeclareMathOperator{\dcl}{dcl} \DeclareMathOperator{\gr}{Gr}
\DeclareMathOperator{\gl}{GL} \DeclareMathOperator{\agl}{AGL}
\DeclareMathOperator{\aut}{Aut} \DeclareMathOperator{\id}{Id}
\DeclareMathOperator{\stab}{Stab} \DeclareMathOperator{\fr}{Fr}
\DeclareMathOperator{\lie}{Lie} \DeclareMathOperator{\inn}{Inn}
\DeclareMathOperator{\zd}{dim} \DeclareMathOperator{\fix}{Fix}
\DeclareMathOperator{\im}{im} \DeclareMathOperator{\dom}{dom}
\DeclareMathOperator{\gal}{Gal} \DeclareMathOperator{\ch}{char}
\DeclareMathOperator{\reg}{reg} \DeclareMathOperator{\fin}{fin}
\DeclareMathOperator{\td}{trdeg} \DeclareMathOperator{\exx}{ex}
\DeclareMathOperator{\ld}{lindeg} \DeclareMathOperator{\theo}{Th}
\DeclareMathOperator{\hght}{height}\DeclareMathOperator{\df}{def}
\DeclareMathOperator{\eq}{eq}
\DeclareMathOperator{\card}{Card}\DeclareMathOperator{\arc}{Arc}
\DeclareMathOperator{\var}{Var}
\DeclareMathOperator{\tp}{tp}\DeclareMathOperator{\stp}{stp}\DeclareMathOperator{\dd}{d}\DeclareMathOperator{\gra}{graph}
\DeclareMathOperator{\trd}{trdeg}\DeclareMathOperator{\loc}{locus}\DeclareMathOperator{\en}{End}
\DeclareMathOperator{\mr}{RM}\DeclareMathOperator{\fron}{fr}\DeclareMathOperator{\grap}{graph}
\DeclareMathOperator{\mlt}{mlt}\DeclareMathOperator{\cl}{cl}\DeclareMathOperator{\omin}{omin}
\DeclareMathOperator{\inte}{int}
\DeclareMathOperator{\rk}{rk}\DeclareMathOperator{\vect}{vect}
\DeclareMathOperator{\sgn}{sgn} \DeclareMathOperator{\st}{st}
\DeclareMathOperator{\rv}{rv}

%\setcounter{section}{0}
\newtheorem{introtheorem}{Theorem}
\newtheorem*{nthm}{Theorem}
\newtheorem{ttheorem}{Theorem}
\newtheorem{pproblem}{Problem}
\newtheorem{theorem}{Theorem}[section]
\newtheorem{acknowledgment}[theorem]{Acknowledgment}
\newtheorem{algorithm}[theorem]{Algorithm}
\newtheorem{axiom}[theorem]{Axiom}
\newtheorem{case}[theorem]{Case}
\newtheorem{claim}{Claim}[theorem]
\newtheorem{conclusion}[theorem]{Conclusion}
\newtheorem{condition}[theorem]{Condition}
\newtheorem{conjecture}[theorem]{Conjecture}
\newtheorem{corollary}[theorem]{Corollary}
\newtheorem{criterion}[theorem]{Criterion}
%\newtheorem{definition}[theorem]{Definition}
%\newtheorem{example}[theorem]{Example}
\newtheorem{exercise}[theorem]{Exercise}
\newtheorem{fact}[theorem]{Fact}
\newtheorem{lemma}[theorem]{Lemma}
\newtheorem{problem}[theorem]{Problem}
\newtheorem{proposition}[theorem]{Proposition}
%\newtheorem{remark}[theorem]{Remark}
\newtheorem{solution}[theorem]{Solution}
\newtheorem{summary}[theorem]{Summary}
\newtheorem*{assumption}{Assumption}

%Axioms:

\newcommand{\efi}{\hyperlink{Aefbi}{(EfI)}}
\newtheorem*{efbi}{\fbox{{\large A}} \hypertarget{Aefbi}{EfI}}


\newcommand{\gendif}{\hyperlink{Agen-dif}{(Gen-Dif)}}
\newtheorem*{gen-dif}{\fbox{{\large A}} \hypertarget{Agen-dif}{Gen-Dif}}


\newcommand{\VJP}{\hyperlink{Avjp}{(VJP)}}
\newtheorem*{vjp}{\fbox{{\large A}} \hypertarget{Avjp}{VJP}}

\newcommand{\MVTay}{\hyperlink{Amvtay}{(MVTay)}}
\newtheorem*{mvtay}{\fbox{{\large A}} \hypertarget{Amvtay}{MVTay}}

\newcommand{\BCen}{\hyperlink{Abcen}{(B-Cen)}}
\newtheorem*{bcen}{\fbox{{\large A}} \hypertarget{Abcen}{B-Cen}}

\newcommand{\minballn}{\hyperlink{Amin-ball}{(Cballs)}}
\newtheorem*{min-balln}{\fbox{{\large A}} \hypertarget{Amin-ball}{Cballs}}

%\newcommand{\Tay}{(Tay)}

\theoremstyle{definition}
%\newtheorem{introtheorem}{Theorem}
\newtheorem{definition}[theorem]{Definition}
\newtheorem{example}[theorem]{Example}
\newtheorem{remark}[theorem]{Remark}
\newtheorem{question}[theorem]{Question}
\newtheorem{notation}[theorem]{Notation}

% BlackBoard Shortcuts:

\newcommand{\Cc}{{\mathbb{C}}}
\newcommand{\Rr}{{\mathbb{R}}}
\newcommand{\Nn}{{\mathbb{N}}}
\newcommand{\Qq}{{\mathbb{Q}}}
\newcommand{\Kk}{{\mathbb{K}}}
\newcommand{\Pp}{{\mathbb{P}}}
\newcommand{\Zz}{{\mathbb {Z}}}

\newcommand{\m}{\textbf{m}}
\newcommand{\bk}{\textbf{k}}

%Mathcal Shortcuts

\newcommand{\CB}{{\mathcal B}}
\newcommand{\CD}{{\mathcal D}}
\newcommand{\CL}{{\mathcal L}}
\newcommand{\CK}{{\mathcal K}}
\newcommand{\CN}{{\mathcal N}}
\newcommand{\CH}{{\mathcal H}}
\newcommand{\CR}{{\mathcal R}}
\newcommand{\CM}{{\mathcal M}}
\newcommand{\CE}{{\mathcal E}}
\newcommand{\CC}{{\mathcal C}}
\newcommand{\CT}{\mathcal T}
\newcommand{\CO}{{\mathcal O}}
\newcommand{\CI}{\mathcal{I}}
\newcommand{\CF}{{\mathscr F}}
\newcommand{\CG}{{\mathcal G}}
\newcommand{\CS}{\mathcal S}
\newcommand{\CP}{\mathcal P}
\newcommand{\CZ}{\mathcal Z}
\newcommand{\CU}{\mathcal U}

\newcommand{\x}{\bar x}
\newcommand{\y}{\bar y}
\newcommand{\0}{\emptyset}
\newcommand{\vphi}{\varphi}
\newcommand{\rest}{\upharpoonright}

%\newcommand{\ax6}{Axiom 6}


%\def\uth{\text{U}^{\text{\th}} }
\newcommand{\fgeq}{\succcurlyeq}
\newcommand{\fleq}{\preccurlyeq}
\renewcommand{\phi}{\varphi}
\renewcommand{\a}{\bar a}
\renewcommand{\b}{\bar b}
\renewcommand{\x}{\bar x}
\renewcommand{\y}{\bar y}

\newcommand{\RV}{\mathrm{RV}}
\newcommand{\ad}{\mathrm{Ad}}

\def\thind{\mathop{\mathpalette\Ind{}}^{\text{\th}} }
\def\nthind{\mathop{\mathpalette\Notind{}}^{\text{\th}} }
\def\uth{\text{U}^{\text{\th}} }
\def\ur{\text{U}}
\def\tho{\text{\th}}
\def\po{\le_p}
\def\bm{\mathfrak m}
\def\la{\langle}
\def\ra{\rangle}
\def\qp{\mathbb Q_p}
\def\zp{\mathbb Z_p}
\def\oD{\overline{D}}
\def\dpr{\mathrm{dp\text{-}rk}}
\def\sub{\subseteq}
\def\drk{\mbox{D-rk}}
\def\adrk{\mbox{aD-rk}}

\DeclareMathOperator{\Tho}{\tho}


\newenvironment{claimproof}[1][\proofname]
{%
	\proof[#1]%
	\renewcommand*\qedsymbol{$\square$ (claim)}%
}
{%
	\endproof%
}

\newenvironment{subclaimproof}[1][\proofname]
{%
	\proof[#1]%
	\renewcommand*\qedsymbol{$\square$ (subclaim)}%
}
{%
	\endproof%
}

\title{A note on uniform finiteness in weakly o-minimal theories}

\author{Yatir Halevi}
\address{Faculty of Mathematics\\Technion - Israel Institute of Technology\\Haifa\\ Israel.}
 \email{yatirh@math.technion.ac.il}

\author{Assaf Hasson}
\address{Department of Mathematics, Ben Gurion University of the Negev, Be'er-Sheva, Israel}
\email{hassonas@math.bgu.ac.il}

\author{Ya'acov Peterzil}
\address{Faculty of Natural Sciences, Department of Mathematics, University of Haifa, Haifa, Israel}
\email{kobi@math.haifa.ac.il}

\date{\today}
%\maketitle
\oddsidemargin .8cm
\evensidemargin .8cm
\begin{document}

 \thanks{The first and second authors were partially supported by ISF grant No. 555/21. The third author was supported by ISF grant No. 290/19.}
    
	\maketitle
     \begin{abstract}
         For an $\aleph_0$-saturated weakly o-minimal expansion of an ordered group $\CM$, it is shown that $\CM^{\eq}$ has uniform finiteness if and only if the collection of definable convex subgroups of $\CM$ has uniform finiteness. If $\CM$ expands an ordered field, then considering definable convex valuation subrings is sufficient. The results use a criterion of Johnson for uniform finiteness in $\CM^{\eq}$, \cite{JohnUF}.
        
        % It is also shown that given a pure real closed valued field, expanding the value group and residue field while retaining their weak o-minimality preserves weak o-minimality of the valued field. As a consequence, 
        
        In addition, it is shown that uniform finiteness in $\CM^{\eq}$ may fail for weakly o-minimal expansions of fields.


        % In the appendix an example of a weakly o-minimal expansion of a real closed field $\CK$ where uniform finiteness fails in $\CK^{\eq}$ is given.
     \end{abstract}
    \section{Introduction}

A definable (or interpretable) set $X$ in a first-order structure $\CM$ has \emph{uniform finiteness} if for any definable family of subsets $\{D_t\subseteq X: t\in T\}$ there is a natural number $n\in \mathbb{N}$ for which $|D_t|=\infty\iff |D_t|>n$. A structure $\CM$ has uniform finiteness if all its definable subsets have uniform finiteness.  Many well-behaved structures have uniform finiteness; but not all, for example $(\mathbb{Z},+,<)$.

Uniform finiteness in $\CM^{\eq}$ is a different matter (see Definition \ref{D: uniform} for a precise meaning of the term): The p-adic field $\mathbb{Q}_p$ has uniform finiteness, but its interpretable value group does not.

In the present note we study uniform finiteness in weakly o-minimal theories. Recall that a linearly ordered structure $\CM$ is \textit{weakly o-minimal } if every definable subset of $M$ is a finite union of convex sets. Throughout we will be working only with $\aleph_0$-saturated weakly o-minimal structures, so we assume that the theory  of the structure is  weakly o-minimal. Weakly o-minimal theories, essentially by definition, have uniform finiteness. So we focus on uniform  finiteness in imaginary sorts. Our main results are: 

\begin{introtheorem}[Theorems \ref{T: convex subgroups is enough} and \ref{T:Main, UF iff conv val rings}]
    Let $\CM$ be an $\aleph_0$-saturated weakly o-minimal expansion of an ordered group. Then 
    \begin{enumerate}
        \item $\CM^{eq}$ has uniform finiteness if and only if the collection of definable convex\footnote{By \cite[Lemma 5.2]{MacMaSt} the convexity assumption is redundant. We keep it here and throughout for improved readability.}  subgroups of $\CM$ has uniform finiteness. 
        \item If, moreover,  $\CM$ is an expansion of an ordered field, then $\CM^{eq}$ has uniform finiteness if and only if the collection of  definable convex valuation subrings of $\CM$ has uniform finiteness.  
    \end{enumerate}
\end{introtheorem}
Our proof builds on a general criterion of Johnson, \cite[Theorem 2.3]{JohnUF}, for a structure $\CM$ to have uniform finiteness in all imaginary sorts. Specifically, Johnson proves that $\CM^{\eq}$ has uniform finiteness if and only if every definable set of \emph{unary imaginaries} has uniform finiteness. A definable set $X$ in $\CM^{\eq}$ is a {\em set of unary imaginaries} if $X$ is a collection (of codes)  of pairwise distinct subsets of $M$.

Using this result, we also prove an Ax-Kochen-Ershov style transfer result, namely that for a weakly o-minimal ordered field $\CM$, $\CM^{\eq}$ has uniform finiteness if and only if for some (equivalently, any) definable valuation $v$ (possibly trivial) the collections of all $M$-definable convex subgroups of $\Gamma_v$ (the value group) and of the additive group of $\bk_v$ (the residue field) have uniform finiteness (see Corollary \ref{C: UF for wom valued fields} for a more precise statement). This result extends our proof from \cite{HaHaPeGps}, asserting that power bounded T-convex expansions of o-minimal fields have uniform finiteness. 

To show that the above results are not vacuous, we give in the appendix an example of a weakly o-minimal expansion of a real closed valued field $\CK$ where uniform finiteness fails in $\CK^{\eq}$. It is a consequence of the following theorem together with  \cite[Example A.6]{HaHaPeSs}.

\begin{introtheorem}[Theorem \ref{T:appendix theorem}]
   Let $\CK$ be a pure real closed valued field. Let $\tilde k$, $\tilde \Gamma$ be weakly o-minimal expansions of the $\CK$-induced structure on the residue field and value group sorts. Then the resulting expansion  $\widetilde{\CK}$ of $\CK$ is weakly o-minimal. 
   % be structure $\CK$ together with weakly o-minimal expansions of of $\bk$ and $\Gamma$. Then $\widetilde{K}$ is weakly o-minimal as well.
\end{introtheorem}

\subsection*{Conventions} Throughout, given a structure $\CM$, unless specified otherwise, by a definable set we mean a definable set in $\CM^{\eq}$ with parameters.

\section{Definable cuts in ordered abelian groups}

The first set of definitions and facts are mostly known and taken, for example, from  \cite{FVKNar,Tressl}. We repeat them for the reader's convenience, and in order to fix some notation. 


Let $(M,<,+)$ be a linearly ordered group. By {\em a cut in $M$} we mean $\mathcal C=(C_1,C_2)$, where $C_1$ is a non-trivial initial segment of $M$ and $C_2$ is its complement. 
Abusing terminology, we will identify the cut $\CC$ and the corresponding initial segment $C_1$.

\begin{definition}\label{D: Stab}
Given a cut $\CC=(C_1,C_2)$ in $M$, the stabilizer group of $\CC$, also known as \emph{the invariance group of $\CC$} is  $G(C_1)=G(\CC)=\{\delta\in M: \delta+C_1=C_1\}$. 
\end{definition}

The invariance group $G(\CC)$ is a convex (possibly trivial) subgroup of $(M,+)$. The following is standard:

\begin{definition}\label{D: cuts}
	Let $\CC=(C_1,C_2)$ be a cut. 
	\begin{enumerate}
		\item  $\CC$ is {\em valuational} if $G(\CC)\neq \{0\}$ and {\em non-valuational}, otherwise.
		\item A non-valuational cut in $M$ is \emph{rational} if it is realized (i.e., $C_1$ has a supremum) in $M$, and \emph{irrational} otherwise. % A Non-valuational cut  $\CC=(C_1,C_2)$ are either {\em rational} or {\em irrational}, depending on whether $C_1$ has a supremum in $M$ or not.
		\item We say that a non-empty  initial segment  of $M$ is (non-)valuational or (ir)rational if its corresponding cut is such.
	\end{enumerate}	
\end{definition}


   If $\CC$ is a valuational cut in $M$ expanding a real closed field, then $\{x: xG\sub G\}$ is a non-trivial valuation ring, \cite[Lemma 6.2]{MacMaSt}. 
   

For a convex set $D$ we write, as usual, $x<D$ to mean $x<c$ for all $c\in D$.  We also write $x\le D$ for $x<D\vee x\in D$. With this notation the following is immediate: 

\begin{lemma}\label{L:all cuts}
	Let $\CC=(C_1,C_2)$ be  a cut and  $H:=G(C)$. Then exactly one of the following holds:
\begin{enumerate}
    \item[(i)] $\CC$ is  irrational and non-valuational.
    \item[(ii)] There is a coset $S$ of $H$, possibly a singleton, such that $C_1=\{x\in M:x\,\Box\, S\}$, with $\Box\in \{<,\leq\}$. 
    \item[(iii)] $\CC$ is valuational  and for $\pi:M\to M/H$ the natural quotient map, $\pi(C_1)$ is an irrational non-valuational initial segment. 
\end{enumerate}	
\end{lemma}
 

\begin{definition}
	A valuational cut $\CC=(C_1,C_2)$ is called \textit{of type (iii)}, if $\pi(C_1)$ is irrational non-valuational in the quotient group $M/G(C_1)$.
\end{definition}

\begin{remark}
    The literature contains finer classification of cuts, for example see \cite[Figure 1]{FVKNar}. In the terminology of \cite{FVKNar}, type (iii) cuts are non-ball cuts with non-trivial invariance group. In the terminology of \cite{Tressl}, type (iii) cuts are cuts with sign $0$ and non-trivial invariance group.
\end{remark}

\begin{example}
    Consider $M=(\mathbb{Q},+)^2$ with the lexicographic ordering. The cut $\mathcal{C}$ induced by $(\sqrt{2},0)$ is of type (iii). Indeed, $G(\mathcal{C})=0\times\mathbb{Q}$ and the induced cut on $M/G(\mathcal{C})$ is irrational (see \cite[Ex3]{FVKNar}).

    An example of a cut of type (iii) in a weakly o-minimal expansion of a field can be found in \cite[Example 3.11(C)]{Tressl-modcomp}.
\end{example}


Since every definable subset of a weakly o-minimal structure is a finite boolean combination of definable initial segments we conclude:

\begin{lemma}\label{L: definable sets in wom}
	If $\CM=(M,+,<,\dots)$ is a weakly o-minimal expansion of an ordered group then every definable subset of $M$ is a boolean combination of definable initial segments as in Lemma \ref{L:all cuts} (i)-(iii).
\end{lemma}
 

The next lemma will not be used in the sequel, but may be of independent interest. It is a generalization of \cite[Lemma 3.1]{HaHaPeSs}. For a convexly valued ordered field $\CK$ we say that the value group $\Gamma$ (resp. the residue field $\bk$) is \emph{definably complete} if any $\CK$-bounded definable subset of $\Gamma$ (resp. $\bk$) has a supremum. With this terminology we have: 

\begin{lemma}\label{L: o-minimal group and res}
	Assume that $(\CK,+,\cdot,<,\CO,\dots)$ is an expansion of a convexly valued ordered field, such that 
	$\bk$ and $\Gamma$ are definably complete.
   
	Every non-trivial definable convex subgroup of $(K,+)$ is a ball.  
    
    For any valuational cut $\CC=(C_1,C_2)$ in $K$ of type (iii), $G(\CC)$ is a closed ball.
\end{lemma} 
\begin{proof}
    Let $H$ be a non-trivial definable convex subgroup of $(K,+)$. Consider balls of the form $B_{\geq r}(0)$ and $B_{>r}(0)$ for $r\in \Gamma$. As $H$ is a convex subgroup it is comparable by inclusion to every ball (open or closed).  If $B_{\geq r}(0)\subseteq H$ for every $r\in \Gamma$ then $H=K$, so the set  $\{r:B_{\geq r}(0)\sub H\}$ must be bounded below. By definable completeness of $\Gamma$, it has an infimum $r_0$.
    If $r_0$ is not a  minimum then it is easy to see that $H=B_{>r_0}(0)$ an open ball, so assume $r_0$ is a minimum. Namely, $B_{>r_0}(0)\subsetneq H\subseteq B_{\geq r_0}(0)$.

     As $B_{\ge r_0}(0)/ B_{> r_0}(0)\cong \bk$, $H$ is mapped via this isomorphism onto a non-zero convex subgroup of $(\bk, +)$, which by the definable completeness of $\bk$ must equal to $\bk$, hence   $H=B_{\geq r_0}(0)$. Thus, $H$ is either a closed or an open ball.

	Now assume that $C_1$ is a valuational initial segment of type (iii). So $H=G(C_1)$ is nontrivial convex and $\pi(C_1)\subseteq K/H$ is irrational and non-valuational. After rescaling we may assume that $H$ is a ball of radius $0$. We claim that $H=\CO$. 
    
    Indeed, if $H=\m$, the maximal ideal, then there exists $z\in \CO^\times$ which does not lie in $H$, i.e. there are elements $x\in C_1$ and $y\notin C_1$ such that $v(y-x)=0$. So there is a closed ball of radius $0$  containing both $x$ and $y$. Translating, we may assume that this closed ball is $\CO$. Thus, $\pi:K\to K/\m$ sends $\CO$ onto $\bk$ and sends $C_1\cap \CO$ to a nontrivial cut in $\bk$. By definable completeness of $\bk$, $\pi(C_1)$ has a supremum, contradicting the assumption that it is irrational.
\end{proof}    

Keeping the notation of the lemma, if $\CK$ is weakly o-minimal, so are $\Gamma$ and $\bk$. In this setting, $\Gamma$ and $\bk$ are definably complete if and only if they are o-minimal (and therefore stably embedded, see \cite{HaOn2a}).

 Below, we say that a cut $\CC$ is called {\em a ball-cut} if there exists a ball $B\sub K$ such that $C=\{x\in K: x\, \square\, B\}$, where $\square \in \{<,\leq \}$.
\footnote{In \cite{FVKNar} the notion of ball-cut is defined similarly, allowing $B$ to be any convex subgroup of $(K,+)$.}
 By what we have seen thus far, we may conclude:
\begin{corollary}
	If $\CK$ is a weakly o-minimal expansion of a valued field and both $\Gamma$ and $\bk$ are o-minimal then every definable subset of $K$ is a boolean combination of non-valuational cuts, ball-cuts and valuational cuts of type (iii), with stabilizers a closed ball.
\end{corollary}


\section{Uniform finiteness in $\CM^{\eq}$}
 In this section we prove the main theorem stated in the introduction. For that, we first refine the definition of uniform finiteness from the introduction; it is based on Johnson's  \cite{JohnUF}. 

\begin{definition}\label{D: uniform}
    Let $\CM$ be a sufficiently saturated structure.
    \begin{enumerate}
        \item A definable (by our convention, possibly interpretable) set $X$ \emph{has uniform finiteness} if for every definable family $\{Y_t\subseteq X:t\in T\}$ there is a natural number $n\in \mathbb{N}$ such that for all $t\in T$: $|Y_t|=\infty \iff |Y_t|>n.$
        \item Let $\CD$ be a (not necessarily definable) collection of definable subsets of $\bigcup_n M^n$. A definable set $X$ is \emph{a set of imaginaries from $\CD$} if there is a definable relation $R\subseteq X\times M^n$ such that
        \begin{enumerate}
            \item $x\mapsto R_x:=\{m\in M^n:(x,m)\in R\}$ is injective, and
            \item for every $x\in X$, $R_x\in \CD$.
    \end{enumerate}
    \item 
          We say that \emph{$\CD$ has uniform finiteness} 
         if any definable set of imaginaries from $\CD$ has uniform finiteness.
         \end{enumerate}

\end{definition}


With this definition, Johnson's theorem \cite[Theorem 2.3]{JohnUF} states  that $\CM^{\eq}$ has uniform finiteness if and only if the collection of all definable subsets of $M=M^1$ has uniform finiteness. When $\CM$ is weakly o-minimal we can reduce to an even simpler collection of sets, as discussed below.  \\


\textbf{From now on we fix  $\CM=(M,+,<,\dots)$, a sufficiently saturated expansion of a weakly o-minimal ordered group}. 

\begin{lemma}\label{L: initial segments}
   If the collection of initial segments  of $M$ has uniform finiteness then $\CM^{\eq}$ has uniform finiteness.
\end{lemma}
\begin{proof}
By Johnson's theorem, it suffices to show that the collection of definable subsets of $M$ has uniform finiteness. To each definable $X\sub M$ we let $C(X)$ denote the (finite) collection  of 
initial  segments  determined by its convex components (each convex set determines two initial segments).  For a set of unary codes $Y$ we let $C(Y):=\bigcup\limits_{X\in Y} C(X)$. 

Let us verify that $Y$ is finite if and only if $C(Y)$ is finite. Indeed, a set $C(Y)$ of $n$ initial segments gives rise to at most $f(n)=2^{2^n}$ boolean combinations of these segments, thus there are at most $f(n)$-many definable sets $X\sub M$ whose code belongs to $Y$, namely $|Y|\leq f(n)$. The other direction is obvious. 

This implies that uniform finiteness of the collection of initial segments implies uniform finiteness of the collection of definable subsets of $M$, as claimed. 
\end{proof}


Our goal here is to establish which  definable sets in $\CM^{eq}$ necessarily have uniform finiteness.
An obvious obstruction  is the existence of a definable infinite discrete linear order. We start by noting that this is, in fact, the only obstacle to $\CM^{eq}$ having uniform finiteness. 

\begin{lemma}\label{L: eliminating infty equiv nondiscrete}
	$\CM^{\eq}$ does not have uniform finiteness if and only if there exists a definable convex equivalence relation $E$ on $M$ such that $(M/E,<)$ infinite and discrete.
\end{lemma}
\begin{proof}
	Infinite discrete linear orders fail uniform finiteness (as witnessed by the family of intervals). So interpretable discrete linear orders contradict uniform finiteness for $\CM^{eq}$. 
    
Conversely, if $\CM^{\eq}$ does not have uniform finiteness, we prove that $\CM$ interprets an infinite discrete linear order.  By Lemma \ref{L: initial segments}, failure of uniform finiteness appears in some definable family of sets of initial segments of $M$. The initial segments are linearly ordered by inclusion. If $\{Y_t\}_{t\in T}$ is a family of codes such that each $Y_t$ consists of codes of initial segments, then each $Y_t$ is linearly ordered, uniformly in $t$. In particular, if $Y_t$ is finite it is discrete with respect to its natural ordering. 

Thus, if $\{Y_t\}_{t\in T}$ fails uniform finiteness, by compactness and saturation, there is some $t_0\in T$ such that $Y_{t_0}$ is infinite and discretely ordered. Given such a set $Y_{t_0}$, denote for $x\in M$ the set $C(x):=\{y\in Y_{t_0}: x\in y\}$.  The  equivalence  relation $E(x_1,x_2)$ on $M$ given by $C(x_1)=C(x_2)$ has the desired properties 
\end{proof}

\subsection{Uniform finiteness in groups}


Lemma \ref{L: initial segments} reduces uniform finiteness in $\CM^{\eq}$ to uniform finiteness of the collection of all initial segments of $M$. The aim of this section is to reduce it further to the collection of all convex subgroups. We first prove this important observation.

\begin{lemma}\label{L: eliminating non-valuational}
The collection of all non-valuational (possibly rational) definable cuts has uniform finiteness. 
\end{lemma}
\begin{proof} 
Assume that $\{Y_t:t\in T\}$ is a definable family of sets of unary codes, each consisting of (codes of) non-valuational initial segments.  For each $t\in T$ denote $I_t:=\{y\in Y_t: y \text{ is a code of a rational cut}\}$ and $J_t$ its complement in $Y_t$. By uniform finiteness in $M$ there is $n\in \Nn$ such that $|I_t|<n$ if and only if $|I_t|<\infty$. So we may assume that $Y_t=J_t$ for all $t$. 


As in the proof of Lemma \ref{L: eliminating infty equiv nondiscrete}, compactness and saturation imply that if there are finite $Y_t$ of unbounded size, then there is $t_0$ such that $Y_{t_0}$ is discrete with respect to inclusion. Let us see that this contradicts weak o-minimality. 

Let $\{y_i\}_{i\in \Nn}\sub Y_{t_0}$ be such that for all $i\in \Nn$ $y_{{i+1}}$ is an immediate successor in $Y_{t_0}$ of $y_{i}$.  By compactness, there is $\alpha>0$ in $M$,  such that for all $i$ there are $y_i<x_1<x_2<y_{i+1}$ with $x_2-x_1>\alpha$. 
%$y_{i+1}-y_{i}>\alpha>0$ for all $i\in \Nn$.
For $y\in Y_{t_0}$ let $s(y)$ denote the successor of $y$ in $Y_{t_0}$. By our choice of $\alpha$, the set of  $y\in Y_{t_0}$, such that  there are $y<x_1<x_2<s(y)$ with $x_2-x_1>\alpha$ is infinite (since it contains $\{y_i\}_{i\in \Nn}$). 

For every code $y\in Y_{t_0}$, let $D_y$ be the convex subset of $M$ given by 
$\{x\in M: x<y< x+\alpha \}$. Since $y$ is non-valuational, $D_y\neq \0$. 
Then $\bigcup\{D_y: y\in Y_{t_0}\}$ is a definable set consisting of infinitely many convex components. This contradicts weak o-minimality.  
\end{proof}



\begin{lemma} \label{L: eliminating type (iii)}
	Assume that the collection of definable convex subgroups of $M$ has uniform finiteness. Then the collection of definable valuational initial segments of type (iii) has uniform finiteness.
\end{lemma}
\begin{proof}
	Assume that we have a definable family $\{Y_t:t\in T\}$, with each $Y_t$ consisting of (codes of) initial segments, $C\sub M$, with stabilizers $H=G(C)$, such that $\pi(C)\sub M/H$ is irrational non-valuational.  

    
	 For each $t\in T$, let $H(t)=\{G(C):C\in Y_t\}$. Then $H(t)$ is a set of convex subgroups of $(K,+)$ and  $|H(t)|\leq |Y_t|$. By the assumption, there is some $n\in \mathbb{N}$, such that $H(t)$ is finite if and only if $|H(t)|\leq n$. Thus, since we want to bound the size of the finite $Y_t$, we may assume that for every $t\in T$, $|H(t)|\leq n$. Since the groups in $H(t)$ are linearly ordered by inclusion, we may partition each $Y_t$ according to the associated group, and assume that for every $t\in T$, $|H(t)|=1$.

     
Assume now that the result fails.
    Then, there exists a strictly increasing sequence of natural numbers $n_k$, and for every $k$, some $t_k\in T$, such that $n_k=|Y_{t_k}|$. Thus,  there is some $t_0$ for which $Y_{t_0}$ is infinite and discretely ordered, and there is some fixed $H$ such that for all $C\in Y_{t_0}$, $G(C)=H$.
    
    Consider now the family of non-valuational initial segments inside $M/H$: $\{\pi(C):C\in Y_{t_0}(H)\}$. Since each such $C$ consists of cosets of $H$, the map $C\mapsto C/H$ is injective and order preserving on $Y_{t_0}(H)$, hence it contains a relatively convex subsequence of order $\omega$.
	
	
	This in turn gives rise to a definable family of  sets of irrational non-valuational cuts in the weakly o-minimal group (so dense) $M/H$, whose finite sets are unbounded in size, contradicting  Lemma \ref{L: eliminating non-valuational}. 
\end{proof}

\begin{lemma}\label{L:UF grp implies UF cosets}
    Assume that the collection of definable convex subgroups of $M$ has uniform finiteness. Then the collection of cosets of definable convex subgroups of $M$ also has uniform finiteness.
\end{lemma}
\begin{proof}
    Let $\{Y_t: t\in T\}$ be a definable family with each  element in $Y_t$ a coset of a definable convex subgroup. We can order cosets of convex subgroups by: $a+D<b+E$ if $D<E$ and if $D=E$ then $a+D<b+D$. Since we want to bound the size of those finite $Y_t$, we may assume that for all $t$, $Y_t$ is discretely ordered with respect to this order.
    
    The definable map, mapping $C\in Y_t$ to $C-C$ maps each $Y_t$ to the set $S_t$ of definable convex subgroups associated to the cosets in $Y_t$. By the assumption, there is some natural number $n$ such that if $|S_t|$ is  finite then $|S_t|<n$. Thus assume that $S_t$ is finite for all $t$.

    By saturation we find some $t_0$ for which $Y_{t_0}$ is infinite and as $|S_{t_0}|$ is finite there is some $H\in S_{t_0}$ with infinitely many cosets as elements of $Y_{t_0}$. This gives a definable, infinite, discrete subset of the weakly o-minimal $M/H$, contradiction.
\end{proof}

\begin{theorem}\label{T: convex subgroups is enough}
    Let $\CM=(M,+,<,\dots)$ be a sufficiently saturated weakly o-minimal expansion of an ordered group. If the collection of all definable convex subgroups of $M$ has uniform finiteness then $\CM^{\eq}$ has uniform finiteness.
\end{theorem}
\begin{proof}
    By Lemma \ref{L: initial segments}, it suffices to show that the collection of initial segments of $M$ has uniform finiteness. 
    
	In Lemma \ref{L: definable sets in wom} we classified the different types of initial segments in $\CM$.  Given a definable family $\{Y_t: t\in T\}$  where each $Y_t$ is a set of (codes of) definable initial segments, we have to show that there is $n\in \Nn$ such that for all $t\in T$ such that $|Y_t|<\infty$ it follows that $|Y_t|< n$. For $t\in T$ partition $Y_t$ (definably, uniformly in $t$) into $Y_t^{(i)}$, $Y_t^{(ii)}$ and $Y_t^{(iii)}$ according to whether $y\in Y_t$ is a code for an initial segment of type (i), (ii) or (iii), respectively. It will suffice to show that there is $n$ such that for all $t$ if any one of $Y_t^{(i)}$, $Y_t^{(ii)}$ and $Y_t^{(iii)}$ is finite, it has size at most $n$. 
    
    For $\{Y_t^{(i)}: t\in T\}$ and $\{Y_t^{(iii)}: t\in T\}$  we use  Lemma \ref{L: eliminating non-valuational} and  Lemma \ref{L: eliminating type (iii)}, respectively. For $\{Y_t^{(ii)}: t\in T\}$,  to each initial segment coded by an element $y\in Y_t^{(ii)}$  corresponds a coset $C(y)$ of a definable convex subgroup of $M$. The map $y\mapsto C(y)$ is at most two-to-one, so the conclusion follows from Lemma \ref{L:UF grp implies UF cosets}. 
    \end{proof}

\begin{remark}\label{R: Thm ok for def wom}
    Although Theorem \ref{T: convex subgroups is enough} is stated for weakly o-minimal structures, the theorem also holds for weakly o-minimal definable sets in sufficiently saturated structures, i.e., definable ordered abelian groups $(D,+,<)$ for which every $\CM$-definable subset of $D$ is a finite union of convex definable sets. Namely, we do not need to assume that $D$ is stably embedded.

    More concretely: Let $D$ be a weakly o-minimal definable ordered abelian group in a sufficiently saturated structure $\CM$. If the collection of all definable convex subgroups of $D$ has uniform finiteness then every $\CM$-definable imaginary sort in $D$ has uniform finiteness.
\end{remark}


\subsection{Uniform finiteness in valued fields}


Let $\CK=(K,+,\cdot,<,\dots)$ be a sufficiently saturated weakly o-minimal expansion of an ordered field.\\

Below, by a discrete linear order we mean one where every non-extremal element has an immediate successor and an immediate predecessor. The following is then straightforward.
\begin{lemma}\label{L:union of discrete}
    Let $(X,<)$ be a definable set with a definable linear order in some sufficiently saturated structure $\CM$ and $\{X_t\subseteq X:t\in T\}$ be a definable family.  Let $(t_n)_{n<\omega}$ be such that $X_{t_n}$ are finite and $|X_{t_n}|\to \infty$.

    For any definable subset $Y\subseteq X$,  there exists $t_0\in T$ such that either $X_{t_0}\cap Y$ or $X_{t_0}\cap (X\setminus Y)$ are infinite and discretely ordered. 
\end{lemma}

\begin{theorem}\label{T:Main, UF iff conv val rings}
    Let $\CK$ be as above. Then $\CK^{\eq}$ has uniform finiteness if and only if the collection of all definable convex valuation subrings of $\CK$ has uniform finiteness.
\end{theorem}
\begin{proof}
Obviously if $\CK^{\eq}$ has uniform finiteness then the collection of all definable convex valuation subrings has uniform finiteness. We prove the reverse implication. We first show the following: \\

\noindent $(\dagger)$ Uniform finiteness of the collection of definable convex  valuation rings implies that for every definable valuation $v$ on $K$, the collection of all definable convex subgroups of $\Gamma_v$ has uniform finiteness.\\


Indeed, by Theorem \ref{T: convex subgroups is enough} (and Remark \ref{R: Thm ok for def wom}) it suffices to show that the collection of definable convex subgroups of $\Gamma$ has uniform finiteness, and since for every convex subgroup of $\Gamma$ there corresponds a unique convex valuation subring of $K$ (the corresponding coarsening of $\CO_v$), the result follows.
    
    To prove the theorem, we use Theorem \ref{T: convex subgroups is enough} again. Assume that $\{Y_t:t\in T\}$ is a definable family, with each $Y_t$ a definable set of (codes of) definable (convex) subgroups of $(K,+)$. The collection of all convex subgroups is linearly ordered by inclusion, and since we want to bound the size of those $Y_t$ which are finite, we may assume that for all $t$, $Y_t$ is discretely ordered by inclusion.  
    % Namely, we assume that every $H\in Y_t$ which is not maximal has a successor and every $H\in Y_t$ which is not minimal has a predecessor. We may also assume that each $Y_t$ has a maximal element.

    We assume towards a contradiction that there is no $n$ which bounds the size of all finite $Y_t$. By saturation there exists $t_0$, such that $Y_{t_0}$ is infinite and by our assumption discretely ordered by inclusion.
	

    
    For each  $H\in \bigcup_{t\in T} Y_t$, let $R(H)=\{a\in K: aH\sub H\}$ be the corresponding (convex) valuation ring. The family $\{R(H):H\in \bigcup_{t\in T}Y_{t}\}$ is a uniformly definable chain (with respect to inclusion) of valuation rings  and therefore its intersection $\CO$ is a definable convex valuation ring. Let $\Gamma$ be the corresponding value group and $v:K^\times\to \Gamma$ the corresponding valuation.
	
	For every convex group $H$, let $v(H)=\{v(a):a\in H\}$, a final segment of $\Gamma$. Each such $v(H)$ either has a minimal element in $\Gamma$ or not. By Lemma \ref{L:union of discrete}, after possibly replacing $Y_{t_0}$, either all the elements of $Y_{t_0}$ have a minimum or none of them have. We claim that in each of the two cases $H\mapsto v(H)$ is injective. Assuming this, we get that $\{v(H):H\in Y_{t_0}\}$ is a discrete family of subsets of $\Gamma$,  contradicting $(\dagger)$ (and Theorem \ref{T: convex subgroups is enough}, Remark \ref{R: Thm ok for def wom}). So we prove injectivity in each of these cases.

    
    
	Assume first that no $v(H)$, $H\in Y_{t_0}$ has a minimum. Note that in that case we have $H=v^{-1}(v(H))$, for assume that $v(a)\in v(H)$ for some $a\in K$, which we may assume to be strictly positive.     Then there is $b\in H$, $b>0$, with  $v(b)<v(a)$. It follows that $a<b$, so by convexity $a\in H$. Thus, the map $H\mapsto v(H)$ is injective.
    
    
    
	
    Now assume that for all $H\in Y_{t_0}$,  $v(H)$ has a minimum, denoted by  $\gamma_H$.
	Note that if $\gamma_H=r$ (and recall $H$ is convex) then $B_{>r}(0)\subsetneq H\sub B_{\geq r}(0)$  and therefore $R(H)\subseteq \CO$. By minimality of $\CO$, $R(H)=\CO$ and hence $H=B_{\geq r}(0)$. Hence $H\mapsto v(H)$ is injective and we are done.
\end{proof}

\begin{remark}  
As in Remark \ref{R: Thm ok for def wom}, the theorem also holds for weakly o-minimal definable fields (which are not necessarily stably embedded). Namely, let $K$ be weakly o-minimal definable field in a sufficiently saturated structure $\CM$. If the collection of all definable convex valuation subrings of $K$ has uniform finiteness then every $\CM$-definable imaginary sort in $K$ has uniform finiteness.
\end{remark}


In \cite{HaHaPeSs} we proved that if $\CK$ is power bounded $T$-convex then $\CK^{\eq}$ has uniform finiteness. The following corollary generalizes this result:  


\begin{corollary}\label{C: UF for wom valued fields}
The following are equivalent for $\CK$ as before:
\begin{enumerate}
        \item $\CK^{\eq}$ has uniform finiteness.
        
        \item For any definable valuation $v$ on $K$ (possibly trivial), the collection of all definable convex subgroups of $\Gamma_v$ and $\bk_v$ has uniform finiteness. 

        \item There exists a definable valuation $v$ on $K$ (possibly trivial) such that the collection of all definable convex subgroups of $\Gamma_v$ and $\bk_v$ has uniform finiteness. 

        \item For any definable valuation $v$ on $K$ (possibly trivial), the collection of all definable convex subgroups of $\Gamma_v$ has uniform finiteness.

        \item For any definable valuation $v$ on $K$ (possibly trivial), the collection of all definable convex subgroups of $\bk_v$ has uniform finiteness.

\end{enumerate}
\end{corollary}
\begin{proof}
    $(1)$ easily implies the rest, (2) easily implies (3) by taking the trivial valuation and similarly (5) implies (1). We show the rest.

    $(3)\implies (1)$. Let $v$ be the given valuation, with valuation ring $\CO_v$. If $v$ is trivial then the collection of definable convex subgroups of $K$ has uniform finiteness so we are done by Theorem \ref{T: convex subgroups is enough}. So assume that $v$ is not trivial. 
    
    We use Theorem \ref{T:Main, UF iff conv val rings} and show that the collection of definable convex valuation subrings has uniform finiteness. Let $\{X_t: t\in T\}$ be a definable family of definable convex valuation subrings and let $(t_n)_n$ be such that $|X_{t_n}|\to\infty$. For any $t\in T$, and $R\in X_t$ either $R\subseteq \CO_v$ or $\CO_v\subseteq R$. Thus we may assume that in addition  either $R\subseteq \CO_v$ for every $R\in X_{t_n}$ and natural number $n$ or the other inclusion for all $R\in X_{t_n}$ and natural number $n$.

    If $R\subseteq \CO_v$ for every such $R$, then the induced valuation rings on $\bk_v$ form a family of definable convex subgroups of $\bk_v$ contradicting  $(3)$. If $\CO_v\subseteq R$ for every such $R$ then the induced definable convex subgroups in $\Gamma_v$ contradict $(3)$.


    $(4)\implies (1)$. This is a repeat of the proof of Theorem \ref{T:Main, UF iff conv val rings}, see $(\dagger)$.
\end{proof}

We still do not know the answer to the following:
\begin{question}    
Let $\CK$ be a T-convex valued field (not necessarily power-bounded). Does $\CK^{\eq}$ have uniform finiteness?
\end{question}

\appendix

\section{An example and a  resplendency theorem for weakly o-minimal structures}
In this section we give an example of a weakly o-minimal expansion of a valued field which does not have uniform finiteness in its imaginary sort. To build such an example, we start with a weakly o-minimal expansion $\Gamma$ of an ordered abelian group, admitting an infinite discrete imaginary sort (see, e.g., \cite[Example A.6]{HaHaPeSs}). We then consider the real closed valued field $\mathbb R((t^\Gamma))$ together with an expansion of RCVF by the extra structure on $\Gamma$, and show that it is weakly o-minimal. The key observation is a weakly o-minimal variant of known resplendence results for henselian valued fields. Explicitly, we show that  weakly o-minimal expansions of the value group and the residue field of a (pure) real closed valued field preserve weak o-minimality of the valued field sort.  

% a pure real closed valued field (whose value group, $\Gamma$ is stably embedded. We add structure to $\Gamma$ in such a way that $\Gamma$ remains stably embedded and weakly o-minimal, but $\Gamma{eq}$ admits a definable .

% We start by reminding some folklore results on adding structure to a stably embedded definable sort. See, for example, \cite[Proposition 2.7]{CodEHaJiRi} for more details.
We first recall: Let $\CM$ be a $|T|^+$-saturated structure in a language $\CL$ and $T=\mathrm{Th}(M)$ its theory. Assume that $D$ is a stably embedded sort in $\CM$. Let $\mathcal{D}$ be $D$ with its $\emptyset$-induced structure and $\widetilde{\mathcal{D}}$ an arbitrary expansion of $\mathcal{D}$. It is folklore (see, e.g., \cite[Proposition 2.7]{CodEHaJiRi} for more details)  that letting $\widetilde L$, $\widetilde{\CM}$  be the corresponding expansions, then $D$ is stably embedded in $\widetilde{M}$ and the $\widetilde{L}$ induced structure on $D$ is $\widetilde{\mathcal{D}}$. Below we apply this without further mention to expansions of the RV-sort of (pure) real closed valued fields. 



Let $\CK=(K,\cdot,+,0,1,\CO)$ be a sufficiently saturated model of RCVF. We denote by $\mathbf{RV}=(\mathrm{RV},\cdot,0,1,\oplus)$ the interpretable RV-sort, augmented by an     element $0$ for which we set $\rv(0)=0$ and the ternary relation \[
\oplus(x,y,z):=(\exists a,b)(\rv(a)=x\land  \rv(b)=y\land \rv(a+b)=z).\] 
% Reminder, $\oplus(x,y,z)$ is the relation $\exists a,b\in K$ such that $\rv(a)=x\wedge \rv(b)=y\wedge \rv(a+b)=z$. 

\begin{fact}\label{F:oplus defined}
    For a fixed $x,y\in \mathrm{RV}$ the following are equivalent:
    \begin{enumerate}
        \item There is a unique $z\in \mathrm{RV}$ such that $\oplus(x,y,z)$.
        \item $x=0$ or $y=0$ or $x\neq \rv(-1)y$.
        \item $x=0$ or $y=0$ or $v(x+y)=\min\{v(x),v(y)\}$.
    \end{enumerate}
\end{fact}
\begin{proof}
    $(1)\iff (3)$ is \cite[Proposition 2.4]{flenner} (and the couple of lines after the proof).

    $(2)\iff (3)$ \cite[Proposition 2.2]{flenner}.
\end{proof}

The above allows us to define a partial function: 
\begin{notation}
        For $x, y\in \mathrm{RV}$ we write $x\oplus y$ for the  unique element $z\in \mathrm{RV}$ satisfying $\oplus(x,y,z)$, if such a $z$ is unique. If such a $z$ is not unique $x\oplus y$ is undefined. 
\end{notation}

It is well known that if $(K,v)$ is a henselian valued field of equi-characteristic $0$  then in the structure $\CK=(K,\mathbf{RV},\rv)$, every formula is equivalent to one without field-quantifiers (i.e., quantifiers ranging over variables from the valued field sort), \cite[Proposition 4.3]{flenner}. As a direct consequence, $\mathbf{RV}$ is stably embedded, see for example \cite[Remark A.10]{Ri:analytic}. These two results hold resplendently, i.e. 
% one allows added structure 
they remain valid under arbitrary expansions of the  $\mathbf{RV}$-sort, see the discussion before \cite[Proposition 4.3]{flenner}. As a consequence:

\begin{fact}\label{F:fieldQE-RV}\cite[Proposition 5.1]{flenner}
    Let $\widetilde{\CK}$ be $\CK$ together with some added structure on $\mathbf{RV}$.     Any definable set $X\subseteq K$ is of the form $\{x\in K: (\rv(x-a_1),\dots, \rv(x-a_n))\in S\}$ for some definable subset $S\subseteq \mathrm{RV}^n$ and $a_1,\dots,a_n\in K$.
\end{fact}

Since $\CK$ is real closed, the order on $K$ is definable, and therefore so is the induced order on $\mathrm{RV}$.  Therefore, there is no harm expanding both structures by these orders. From now we let $\CK=(K,\cdot,+,0,1,<,\CO)$ and $\mathbf{RV}=(\mathrm{RV},\cdot,0,1,<,\oplus)$.

% Recall that we say that a definable ordered sort $(D,<)$ of $\CK$ is weakly o-minimal if every definable subset of $D$ is a union of finitely many convex sets.

  \begin{lemma}\label{L:Rvwom implies K wom}
      Let $\widetilde K$ be $K$ together with a weakly o-minimal enrichment of $\mathbf{RV}$ which we denote by $\widetilde{\mathbf{RV}}$. Then $\widetilde K$ is still weakly o-minimal.
  \end{lemma}
\begin{proof}
    Let $X\subseteq K$ be a definable set in $\widetilde K$, by Fact \ref{F:fieldQE-RV} it is of the form \[\{x\in K: (\rv(x-a_1),\dots, \rv(x-a_n))\in S\},\] for some definable subset $S\subseteq \mathrm{RV}^n$.

    For any $1\leq i\leq n$ let $U_{i}=\{x\in K: \bigwedge_j v(x-a_{i})\geq v(x-a_j)\}$. This is a finite union of convex sets and $K=\bigcup_i U_i$.
    % Note that $K=\bigcup_{i_0} U_{i_0}$, indeed for any $x\in K$ choose $i_0$ such that $v(x-a_{i_0})$ is maximal.
    
    
    Let $D_{i}=\{y\in \mathrm{RV}: (y\oplus \rv(a_{i}-a_1),\dots, y\oplus \rv(a_{i}-a_n))\in S\}$. 
    Let us verify that  
    \[X=\bigcup_{i} U_{i}\cap \{x\in K: \rv(x-a_{i})\in D_{i}\}.\]

    Right to left inclusion. Let $x\in U_i$ with $\rv(x-a_i)\in D_i$. In particular, $\rv(x-a_i)\oplus \rv(a_i-a_j)$ is defined for all $j$ and thus is necessarily equal to $\rv(x-a_j)$. This gives that $(\rv(x-a_1),\dots, \rv(x-a_n))\in S$.
    
    Left to right inclusion. Let $x\in X$, then $x\in U_{i}$ for some $i$. To show that $\rv(x-a_{i})\in D_{i}$ it is enough to show that for all $j$, $\rv(x-a_{i})\oplus \rv(a_i-a_j)$ is defined, for then as above it is equal to $\rv(x-a_j)$ and the result follows since $x\in X.$ 
    
    We show that $\rv(x-a_{i})\oplus \rv(a_i-a_j)$ is defined using Fact \ref{F:oplus defined}. Assume that $\rv(x-a_{i})\oplus \rv(a_i-a_j)$ is undefined for some $j$: necessarily $x- a_{i}\neq 0$, $a_j-a_i\neq 0$ and $\rv(x-a_{i})=\rv(a_{j}-a_i)$. By \cite[Proposition 2.2]{flenner}, $v(a_j-a_i-(x-a_i))>v(x-a_i)$, i.e. $v(x-a_j)>v(x-a_{i})$ contradicting $x\in U_{i}$. 
    
    By weak o-minimality of $\widetilde{\mathbf{RV}}$, $D_{i}$ is a finite union of convex sets, thus $\{x\in K:\rv(x-a_{i})\in D_{i}\}$ is also a finite union of convex sets. We conclude that $X$ is a finite union of convex sets.
\end{proof}


 
    Transferring weak o-minimality from the $\mathrm{RV}$ sort to the valued field is the first step, the second is to transfer from $\Gamma$ and $\bk$ to $\mathrm{RV}$. For that we will need to deal with short exact sequences.
    The following is proved in \cite{AsChGeZi} but taken in this form from \cite[Fact 1.74]{Touchard}. Note that as $A$ is divisible, it is pure in $B$.
    % , applied to the short exact sequence of abelian groups $1\to \bk_{>0}\to \mathrm{RV}_{>0}\to \Gamma\to 0$.  

    \begin{fact}\label{F:relQE in short exact}
        Let $\CM=(0\to A\xrightarrow[]{\iota} B\xrightarrow{v}C\to 0)$ be a short exact sequence of abelian groups, with $A$ divisible, allowing expansion of the group structure on $A$ and $C$. 

        Then $A$ and $C$ are stably embedded in $\CM$ and every definable subset of $B$ is a boolean combination of definable sets of the following forms:
        \begin{enumerate}
            \item $\{x\in B: (v(t_1(x)),\dots v(t_n(x))\in V\}$, where the $t_i$ are terms in the group language and $V$ is a definable subset of $C^n$.
            \item $\{x\in B: (\rho(t_1(x)),\dots \rho(t_n(x))\in U\}$, where the $t_i$ are terms in the group language, $U$ is a definable subset of $A^n$ and $\rho(a)=\iota^{-1}(a)$ if $a\in \iota(A)$ and $0$, otherwise.
        \end{enumerate}
    \end{fact}

    Using this result we prove a general weak o-minimality transfer result.

    
    \begin{proposition}\label{P:wom shortexact}
    Let $\CM=(0\to A\xrightarrow[]{\iota} B\xrightarrow{v}C\to 0)$ be a short exact sequence of abelian groups, with $A$ divisible. 

        Let $\widetilde{\CM}=((B,+,0),\widetilde{A}=(A,+,0<,\dots),\widetilde{C}=(C,+,0,<,\dots),\iota,v)$ be $\CM$ together with weakly o-minimal expansions of $A$ and $C$. Then, 

        \begin{enumerate}
            \item There is a unique (definable) order $<$ on $B$ such that with respect to this order, $\iota$ and $v$ are order preserving homomorphisms, $\iota(A)$ is a convex subgroup and $(B,+,<)$ is an ordered abelian group.
            \item The sort $(B,+,0,<)$ (together with its induced structure) is weakly o-minimal.
        \end{enumerate}
    \end{proposition}
    \begin{proof}
        (1) The order on $B$ defined by $b>0$ if and only if $v(b)>0$ or $v(b)=0$ and $\iota^{-1}(b)>0$, is the unique order on $B$ rendering it an ordered abelian group such that $\iota$ and $v$ are order preserving. 
        
        (2) Applying Fact \ref{F:relQE in short exact}, every definable subset of $B$ is a boolean combination of definable sets of the form 
        \begin{enumerate}
            \item $\{x\in B: (v(a_1+m_1x),\dots, v(a_k+m_kx))\in V\}$ for $a_i\in B$, $m_i$ an integer and $V$ a definable set in $\widetilde{C}$ and
            \item $\{x\in B:(\rho(a_1+m_1x),\dots,\rho(a_k+m_kx))\in U\}$ for $a_i\in B$, $m_i$ an integer and $U$ a definable set $\widetilde{A}$, where $\rho(a)=\iota^{-1}(a)$ if $a\in \iota(A)$ and $0\in A$ otherwise. 
        \end{enumerate}
        
        It suffices to show that each kind of these definable sets is a finite union of convex sets.
        
        Let $X\subseteq B$ be a definable set of the first kind. Written differently, it is the same as $v(x)\in \{y: (v(a_1)+m_1y,\dots, v(a_k)+m_ky)\in V\}$. The latter is a definable subset of $\widetilde{C}$ so a finite union of convex sets by weak o-minimality. As $v$ is order preserving, $X$ is a finite union of convex subsets of $B$.

        Let $X\subseteq B$ be a definable set of the second kind. For each $J\subseteq \{1,\dots,k\}$ let $P_J=\{x\in B: a_i+m_ix\in \iota(A)\iff i\in J\}$. It is a finite covering of $B$ so it is enough to show that $X\cap P_J$ is a finite union of convex sets.

        \begin{claim}
            Each $P_J$ is a finite union of  convex sets.
        \end{claim}
        \begin{claimproof}
            Note that $a_i+m_ix\in \iota(A)\iff v(a_i)+m_iv(x)=0$. Setting $Y=\{y\in C:\bigwedge_{i\in J} m_iy=-v(a_i)\wedge \bigwedge_{i\notin J}m_iy\neq-v(a_i)\}$. This is a finite union of convex sets and since $P_J=\{x\in B:v(x)\in Y\}$ it is a finite union of convex sets as well.      
        \end{claimproof}
        
        % Note that $X\cap P_J=\{x\in RV:(\rho(a_1'x^{m_1'}),\dots,\rho(a_k'x^{m_k'}))\in D\}$, where $a_i'=a_i,\, m_i'=m_i$ for all $i\in J$ and $a_i'=1,\, m_i'=0$ for all $i\notin J$. So we may assume that $a_ix^{m_i}\in \iota(\bk)$ for all $i$.
        
        % % By replacing $a_i$ with $0$ for $i\notin  J$ we may replace $X$ with a definable subset where $a_ix^{m_i}\in \iota(\bk)$ for all $i$.

        If $m_i=0$ for every $i\in J$ then $X\cap P_J=\{x\in P_J: (\rho(a_1'),\dots,\rho(a_k'))\in U\}$, where $a_i'=a_i$ if $i\in J$ and $a'_i=0$ otherwise. So either empty or all of $P_J$, and thus either way a finite union of convex sets.

        As in the proof of the claim, if there exists $i\in J$ with $m_i\neq 0$ then either $P_J=\emptyset$ and we are done or $P_J\subseteq v^{-1}(\gamma)$ non-empty for some $\gamma\in C$. Fix an element $b\in v^{-1}(\gamma)$, and let $c_i\in A$ be such that $\iota(c_i):=a_i+m_ib$.
        \begin{claim}
            $X\cap P_J=b+\iota(Y_J)$, for $Y_J=\{y\in A:(d_1(y),\dots,d_k(y))\in U\}$ where $d_i(y)=c_i+m_iy$ if $i\in J$ and $0$ otherwise.
        \end{claim}
        \begin{claimproof}
            Let $x\in X\cap P_J$ so $x=b+\iota(t)$ for some $t\in A$. We will show that $t\in Y_J$. Indeed, for every $i\in J$, $d_i(t)=c_i+m_it=\iota^{-1}(a_i+m_ib+m_i\iota(t))=\iota^{-1}(a_i+m_ix)=\rho(a_i+m_ix)$. The other direction is similar.
        \end{claimproof}

        As $Y_J$ is definable in $A$, it is a finite union of convex sets and thus so is $X\cap P_J$, as required.
    \end{proof}

        We wish to apply this result to a short exact sequence coming from $\mathrm{RV}, \bk$ and $\Gamma$.

    Let $\mathbf{RV}_{\bk,\Gamma}=((\mathrm{RV},\cdot, 0,1),\bk=(\bk,\cdot,+,0,1),\mathbf{\Gamma}=(\Gamma,+,0,<,\infty),\iota,v)$ be the three sorted structure associated  together with a $0$ element to $\bk$ and $\mathrm{RV}$ and an $\infty$ element to $\Gamma$. 
     Whenever $(K,v)$ is a henselian equi-characteristic $0$ valued field, both $\Gamma$ and $\bk$ are stably embedded inside $\mathbf{RV}_{\bk,\Gamma}$, as a pure ordered abelian group and as a pure field, respectively. As noted in \cite[\S 3]{ChSi:inp} the ternary relation $\oplus$ on $\mathrm{RV}$ is also definable in $\mathbf{RV}_{\bk, \Gamma}$.
    
    
    By Proposition \ref{P:wom shortexact}(1) the induced order on $\mathrm{RV}$ is definable.

        
    \begin{lemma}\label{L:wom shortexact of RV}
        Let $\widetilde{\mathbf{RV}_{\bk,\Gamma}}=((\mathrm{RV},\cdot, 0,1,<),\widetilde{\bk}=(\bk,\cdot,+,0,1,<,\dots),\widetilde{\mathbf{\Gamma}}=(\Gamma,+,0,<,\infty,\dots),\iota,v)$ be $\mathbf{RV}_{\bk,\Gamma}$ together with an enrichment of $\bk$ and $\Gamma$ still rendering them weakly o-minimal. Then the $\mathrm{RV}$ sort of $\widetilde{\mathbf{RV}_{\bk,\Gamma}}$ is  weakly o-minimal.
    \end{lemma}
    \begin{proof}
       Let $\mathrm{RV}_{>0}$ be $\rv(K_{>0})$. It is enough to show that $\mathrm{RV}_{>0}$ is weakly o-minimal (there is a definable bijection between $\mathrm{RV}_{>0}$ and $\mathrm{RV}_{<0}$).      Note that $\mathrm{RV}_{>0}$ is a subgroup of $\mathrm{RV}$. Also, $\bk_{>0}$ is a multiplicative subgroup and $\mathrm{RV}_{>0}/\bk_{>0}\cong \Gamma$.

        
         We have a short exact sequence $1\to \bk_{>0}\to \mathrm{RV}_{>0}\to \Gamma\to 0$, $\bk_{>0}$ is divisible because $\bk$ is real closed. 
         
        We are now in the situation of Proposition \ref{P:wom shortexact} where on $\Gamma$ we take the reverse order in order to make $v$ order preserving.
        \end{proof}

    
    
    % \begin{corollary}\label{C:RelQE for RVGammak}
    %     Let $((\mathrm{RV}_{>0},\cdot, 1,<),(\bk_{>0},\cdot,+,1,<,\dots),(\Gamma,+,0,<,\dots),\iota,v)$ be the short exact sequence together with an enrichment of $\bk_{>0}$ and $\Gamma$.

    %     Every definable subset of $\mathrm{RV}_{>0}$ is a boolean combination of definable sets of the form
    %     \begin{enumerate}
    %         \item $\{x\in \mathrm{RV}_{>0}: (v(t_1(x)),\dots v(t_n(x))\in Y\}$, where the $t_i$ are terms in the group language and $Y$ is a definable subset of $\Gamma^n$.
    %         \item $\{x\in \mathrm{RV}_{>0}: (\rho(t_1(x)),\dots \rho(t_n(x))\in Z\}$, where the $t_i$ are terms in the group language, $Z$ is a definable subset of $\bk_{>0}^n$ and $\rho(a)=\iota^{-1}(a)$ if $a\in \iota(\bk_{>0})$ and $1$, otherwise.
    %     \end{enumerate}
    % \end{corollary}
    % \begin{proof}
    %     We have a short exact sequence $1\to \bk_{>0}\to \mathrm{RV}_{>0}\to \Gamma\to 0$, $\bk_{>0}$ is pure in $\mathrm{RV}_{>0}$ and is divisible because $\bk$ is real closed. The only thing yet to verify is that $>$ is definable on $\mathrm{RV}_{>0}$. Indeed, $x>y$ if and only if $v(x)<v(y)$ or $v(x)=v(y)$ and there exists an element $r\in \bk_{>0}$ with $r>1$ and $\iota(r)=xy^{-1}$.
    % \end{proof}
    
    % \begin{lemma}\label{L:wom shortexact}
    %     Let $\widetilde{\mathbf{RV}_{\bk,\Gamma}}=((\mathrm{RV},\cdot, 0,1,<),\widetilde{\bk}=(\bk,\cdot,+,0,1,<,\dots),\widetilde{\mathbf{\Gamma}}=(\Gamma,+,0,<,\infty,\dots),\iota,v)$ be $\mathbf{RV}_{\bk,\Gamma}$ together with an enrichment of $\bk$ and $\Gamma$ still rendering them weakly o-minimal. Then the $\mathrm{RV}$ sort of $\widetilde{\mathbf{RV}_{\bk,\Gamma}}$ is  weakly o-minimal.
    % \end{lemma}
    % \begin{proof}
    %     It is enough to show that $\mathrm{RV}_{>0}$ is weakly o-minimal (there is a definable bijection between $\mathrm{RV}_{>0}$ and $\mathrm{RV}_{<0}$). The advantage is that now $\bk_{>0}$ is divisible.

    %     Applying Corollary \ref{C:RelQE for RVGammak}, every definable subset of $\mathrm{RV}_{>0}$ is a boolean combination of definable sets of the form 
    %     \begin{enumerate}
    %         \item $\{x\in \mathrm{RV}_{>0}: (v(a_1x^{m_1}),\dots, v(a_kx^{m_k}))\in C\}$ for $a_i\in \mathrm{RV}_{>0}$, $m_i$ an integer and $C$ a definable set in $\widetilde{\mathbf{\Gamma}}$ and
    %         \item $\{x\in \mathrm{RV}_{>0}:(\rho(a_1x^{m_1}),\dots,\rho(a_kx^{m_k}))\in D\}$ for $a_i\in \mathrm{RV}_{>0}$, $m_i$ an integer and $D$ a definable set $\widetilde{\bk}_{>0}$. Where $\rho(a)=\iota^{-1}(a)$ if $a\in \iota(\bk)$ and $1\in\bk$ otherwise. 
    %     \end{enumerate}
        
    %     It suffices to show that each kind of these definable sets is a finite union of convex sets.
        
    %     Let $X\subseteq \mathrm{RV}_{>0}$ be a definable set of the first kind. Written differently, it is the same as $v(x)\in \{y: (v(a_1)+m_1y,\dots, v(a_k)+m_ky)\in C\}$. The latter is a definable subset of $\widetilde{\mathbf{\Gamma}}$ so a finite union of convex sets by weak o-minimaity. 

    %     Let $X\subseteq \mathrm{RV}_{>0}$ be a definable set of the second kind. For each $J\subseteq \{1,\dots,n\}$ let $P_J=\{x\in \mathrm{RV}_{>0}: a_ix^{m_i}\in \iota(\bk_{>0})\iff i\in J\}$. It is a finite covering of $\mathrm{RV}_{>0}$ so enough to show that $X\cap P_J$ is a finite union of convex sets.

    %     \begin{claim}
    %         Each $P_J$ is a finite union of  convex sets.
    %     \end{claim}
    %     \begin{claimproof}
    %         Note that $a_ix^{m_i}\in \iota(\bk_{>0})\iff v(a_i)+m_iv(x)=0$. Setting $Y=\{y\in \Gamma:\bigwedge_{i\in J} m_iy=-v(a_i)\wedge \bigwedge_{i\notin J}m_iy\neq-v(a_i)\}$. This a finite union of convex sets and since $P_j=\{x\in \mathrm{RV}_{>0}:v(x)\in Y\}$ it is finite union of convex sets as well.        
    %     \end{claimproof}
        
    %     % Note that $X\cap P_J=\{x\in RV:(\rho(a_1'x^{m_1'}),\dots,\rho(a_k'x^{m_k'}))\in D\}$, where $a_i'=a_i,\, m_i'=m_i$ for all $i\in J$ and $a_i'=1,\, m_i'=0$ for all $i\notin J$. So we may assume that $a_ix^{m_i}\in \iota(\bk)$ for all $i$.
        
    %     % % By replacing $a_i$ with $0$ for $i\notin  J$ we may replace $X$ with a definable subset where $a_ix^{m_i}\in \iota(\bk)$ for all $i$.

    %     If $m_i=0$ for every $i\in J$ then $X\cap P_J=\{x\in P_J: (\rho(a_1'),\dots,\rho(a_k'))\in D\}$, where $a_i'=a_i$ if $i\in J$ and $a'_i=1$ otherwise. So either empty or all of $P_J$, and thus either way a finite union of convex sets.

    %     As in the proof of the claim, if there exists $i\in J$ with $m_i\neq 0$ then either $P_J=\emptyset$ and we are done or $P_J\subseteq v^{-1}(\gamma)$ non-empty for some $\gamma\in \Gamma$. Fix an element $b\in v^{-1}(\gamma)$, and let $c_i\in\bk_{>0}$ be such that $\iota(c_i):=a_ib^{m_i}$.
    %     \begin{claim}
    %         $X\cap P_J=b\cdot \iota(Y_J)$, for $Y_J=\{y\in \bk_{>0}:(d_1(y),\dots,d_k(y))\in D\}$ where $d_i(y)=c_iy^{m_i}$ if $i\in J$ and $1$ otherwise.
    %     \end{claim}
    %     \begin{claimproof}
    %         Let $x\in X\cap P_J$ so $x=b\iota(t)$ for some $t\in \bk_{>0}$. We will show that $t\in Y_J$. Indeed, for every $i\in J$, $d_i(t)=c_it^{m_i}=\iota^{-1}(a_ib^{m_i}\iota(t)^{m_i})=\iota^{-1}(a_ix^{m_i})=\rho(a_ix^{m_i})$. The other direction is similar.
    %     \end{claimproof}

    %     As $Y_J$ is definable in $\widetilde \bk_{>0}$, it is a finite union of convex sets thus so is $X\cap P_J$, as required.
    % \end{proof}

    \begin{theorem}\label{T:appendix theorem}
        Let $\CK$ be a pure real closed valued field and let $\widetilde{K}$ be  the structure $\CK$ together with weakly o-minimal expansions of $\bk$ and $\Gamma$. Then $\widetilde{K}$ is weakly o-minimal as well.
    \end{theorem}
    \begin{proof}
           The $\widetilde{\mathbf{RV}_{\bk,\Gamma}}$ reduct is weakly o-minimal by Lemma \ref{L:wom shortexact of RV}.  As $\oplus$ is definable in this structure (see  \cite[Section 3]{ChSi:inp}), $\widetilde{\mathbf{RV}}$ is a reduct of $\widetilde{\mathbf{RV}_{\bk,\Gamma}}$. By Lemma \ref{L:Rvwom implies K wom}, $\widetilde{K}$ is weakly o-minimal.
    \end{proof}

    
    \begin{corollary}
        There exists a weakly o-minimal valued field $\CK$ for which $\CK^{\eq}$ does not have uniform finiteness.
    \end{corollary}
    \begin{proof}
        In \cite[Example A.6]{HaHaPeSs}, we constructed a weakly o-minimal ordered abelian group $\mathcal{Q}=(Q,+,0,<)$ and a convex equivalence relation $E$ on $Q$ such that $\widetilde{\mathcal{Q}}=(Q,+,0,<,E)$ is still weakly o-minimal but $Q/E$ is discrete. There is no harm in assuming that $\mathcal{Q}$ (and $\widetilde{\mathcal{Q}}$) is $|T|^+$-saturated.
    
        Let $\CK$ be a $|T|^+$-saturated pure real closed valued field with value group $\mathcal{Q}$, e.g. $\mathcal{R}((t^{\mathcal{Q}}))$, for some $\mathbb{R}\prec \mathcal{R}$. By Theorem \ref{T:appendix theorem}, the structure $\widetilde{\mathcal{K}}$ obtained from $\mathcal{K}$ after expanding the value group to $\widetilde{\mathcal{Q}}$ is still weakly o-minimal. As $\widetilde{\mathcal{Q}}^{\eq}$ does not have uniform finiteness, $\widetilde{\mathcal{K}}^{\eq}$ does not have it as well.
    \end{proof}
    
\bibliographystyle{plain}
\bibliography{harvard}

\end{document}
